\respcap{cap:heapav}
\resp{hb-l7} Faça $n$ inserções sucessivas. Inserir é somar 1 em binário, e cada \emph{link} corresponde a um ``vai um''. Contando de 0 até $n$, o ``vai um'' do bit $k$ para o bit $k+1$ ocorre no máximo $\floor{n/2^{k+1}}$ vezes, de modo que o número total de links é $\sum_{k\ge1}\floor{n/2^k}<n$, e o custo total é $O(n)$. Alternativa equivalente: agrupe os elementos em pares formando árvores $B_1$, depois os pares de $B_1$ formando $B_2$, e assim por diante, como numa soma binária de $n$ uns.

\resp{hb-l8} A inserção é o pior caso quando o novo $B_0$ provoca ``vai um'' em todas as posições: o heap contém $B_0,B_1,\dots,B_{k-1}$, isto é, $n=2^k-1$, com todos os bits iguais a 1. As alturas são $0,1,\dots,k-1$ em arestas (ou $1,\dots,k$ em níveis). A inserção faz $k$ links e produz uma única árvore $B_k$.

\resp{hb-l10} A remoção percorre as $k$ raízes, transforma os $h$ filhos da raiz removida no heap $H''=\{B_0,\dots,B_{h-1}\}$ e funde $H''$ com o restante $H'$. O número de links é máximo quando $H'$ também contém $B_0,\dots,B_{h-1}$ e $h$ é o maior possível: $H=\{B_0,B_1,\dots,B_{k-1}\}$ com o máximo na raiz de $B_{k-1}$, isto é, $h=k-1$ (em arestas).

\resp{hb-l11} Falso. $H_1$ com $3=11_2$ nós ($k_1=2$) e $H_2$ com $5=101_2$ nós ($k_2=2$). A fusão tem $8=1000_2$ nós, numa única árvore: $k=1<2$.

\resp{hb-l12} Verdadeiro. A união das listas de raízes tem $k_1+k_2$ árvores e cada link transforma duas árvores em uma; logo $k=k_1+k_2-\#\text{links}\le k_1+k_2$.

\resp{hb-l13} Suponha árvores de posto $k$ formadas por subordinações de árvores de mesmo posto, com cortes controlados por marcas em nós que não são raízes, e meça alturas em arestas. \textbf{Altura mínima}: $\ceil{k/2}$. Pelo lema $\operatorname{grau}(y_i)\ge i-2$, a menor altura possível satisfaz $A(k)=1+A(k-2)$, com $A(0)=0$ e $A(1)=1$. \textbf{Altura máxima}: $k$, atingida pela própria $B_k$; por indução, subordinar duas árvores de posto $h$ e altura $\le h$ dá posto $h+1$ e altura $\le h+1$, e cortes não aumentam a altura. Se as raízes puderem perder filhos livremente, só o \emph{grau} continua sendo $O(\log n)$; a altura pode crescer.

\resp{hb-l14} Ver §\ref{sec:binomial} e §\ref{sec:fib}. Heap binomial: a seleção percorre as raízes (ou mantém um ponteiro); a inserção é a fusão com $\{B_0\}$; a fusão intercala as listas de raízes por ordem e faz links como numa soma binária; a remoção do máximo retira a árvore, transforma os filhos num heap e funde; o aumento sobe trocando com o pai; a diminuição remove o nó (aumentando-o para $+\infty$ e removendo o máximo) e o reinsere; a construção faz $n$ inserções. Heap de Fibonacci: inserção e fusão concatenam listas circulares; a remoção do máximo promove os filhos e consolida por grau; o aumento corta o nó e faz cortes em cascata; a diminuição remove e reinsere.

\resp{hb-r6} (a) $B_0$ é um nó; $B_k$ é obtida ligando duas $B_{k-1}$. Tem $2^k$ nós e a raiz tem grau $k$, que é a ordem da árvore. (b) Como uma lista de raízes de árvores com a condição de heap e \textbf{ordens distintas}; $B_k$ aparece se e somente se o bit $k$ de $n$ é 1. (c) A raiz 10 fica por cima e a raiz 15 vira sua filha; o resultado é uma $B_3$, com 8 nós e raiz de grau 3. (d) O extremo está numa das no máximo $\floor{\log n}+1$ raízes: $O(\log n)$, ou $O(1)$ com ponteiro. (e) $O(\log n)$: a união percorre só as listas de raízes, que têm $O(\log n)$ árvores, fazendo links como numa soma binária; os nós internos não são visitados.

\resp{hb-r9} (a) O heap de Fibonacci é preguiçoso: permite várias árvores de mesmo grau na lista de raízes, consolidando só na remoção, e usa listas circulares duplas e marcas nos nós. (b) Inserção, união, seleção e redução de chave (\emph{decrease-key}) custam $O(1)$ amortizado; \emph{Extract-Min} custa $O(\log n)$ amortizado. (c) A nova chave pode ficar menor que a do pai. Corta-se $x$ com sua subárvore, e ele vai para a lista de raízes, desmarcado. (d) Se o pai já tinha perdido um filho (estava marcado), ele também é cortado, e assim por diante. Isso garante que cada nó perde no máximo um filho antes de ser cortado, o que mantém o grau $O(\log n)$. Com o potencial $\Phi=t+2m$, cada corte em cascata desmarca um nó: o potencial cai 2 e sobe 1, e o corte paga a si mesmo. Por isso o custo amortizado é $O(1)$.

\resp{hb-m1} São $21=10101_2$ e $3=11_2$ elementos. Os links, na ordem, são: $B_0+B_0\to B_1$; esse $B_1$ com o $B_1$ de $H_2$ forma $B_2$; essa $B_2$ com a $B_2$ de $H_1$ forma $B_3$. O resultado é $\{B_3,B_4\}$, isto é, $24=11000_2$.

\resp{hb-m2} (1) $X$ é cortado e vai para a lista de raízes, desmarcado. (2) $Y$, que estava marcado, perdeu o segundo filho: é cortado, desmarcado e vai para as raízes. (3) O mesmo acontece com $Z$. (4) Se o pai $W$ de $Z$ é raiz, nada muda; se não é raiz e está desmarcado, é marcado e a cascata termina; se estava marcado, a cascata continua. Vão para a lista de raízes $X$, $Y$ e $Z$ (e possivelmente outros); ao final, $X,Y,Z$ estão desmarcados e $W$ marcado, se não for raiz.

\resp{hb-m3} Sem as marcas, um nó poderia perder quase todos os filhos sem que nada acontecesse: cortando todos os netos de uma raiz, obtém-se uma árvore de posto $k$ com apenas $k+1$ nós. O tamanho deixaria de crescer exponencialmente com o posto (como $F_{k+2}$), o grau máximo deixaria de ser $O(\log n)$ e a consolidação, que depende do número de graus distintos, ficaria cara.

\resp{hb-t1} \textbf{Heap binomial}. A fusão custa $O(\log n)$ \emph{no pior caso de cada execução} e cresce só logaritmicamente com o número de aeronaves: é previsível. No heap binário, a fusão exige reconstruir o vetor, $\Theta(n)$ por fusão; milhares de fusões com $10^5$ elementos estouram as janelas de tempo. No heap de Fibonacci, a fusão é $O(1)$ apenas \emph{amortizado}: o trabalho adiado aparece de uma vez numa remoção, cuja consolidação pode custar $\Theta(n)$, o que é imprevisível; além disso, a grande vantagem dessa estrutura, o aumento de prioridade em $O(1)$ amortizado, não seria usada, já que as alterações são raras.

\resp{hb-t2} \textbf{Heap binário em vetor}. Os dados ocupam $2\cdot10^6\times16=32\,000\,000$ bytes, cerca de 30,5 MB, abaixo do limite de $41\,943\,040$ bytes. Um heap binomial precisaria de pelo menos 3 ponteiros por nó (pai, filho, irmão), isto é, 24 bytes, mais o grau: pelo menos 40 bytes por nó e $80\cdot10^6$ bytes no total, o que não cabe; o de Fibonacci, com 4 ponteiros, grau e marca, ainda menos. Como não há fusões e as alterações são poucas, inserção e remoção em $O(\log n)$ bastam; o heap binário é o único que cabe e o mais simples.

\resp{hb-t3} \textbf{Heap de Fibonacci}. O aumento de prioridade custa $O(1)$ amortizado, sem retirar e reinserir a tarefa. Uma estimativa do custo total é $5\cdot10^6+8\cdot10^7+5\cdot10^6\log_2(5\cdot10^6)\approx2\cdot10^8$ passos. Nos heaps binário e binomial, cada um dos $8\cdot10^7$ aumentos custa $O(\log n)\approx22$ passos, cerca de $2\cdot10^9$ no total. Como importa o custo total da execução, e não a previsibilidade de cada operação, a garantia amortizada é adequada; as poucas uniões não pesam.

\resp{hb-vetor} Com $k$ elementos e capacidade $c$, vale $c\ge k\ge c/2$ depois da primeira duplicação, logo $\Phi=2k-c\ge0$. Uma inserção que não duplica tem custo 1 e $\Delta\Phi=2$: custo amortizado 3. Uma inserção que duplica (com $k=c$) custa $c+1$, pois copia $c$ elementos e insere um; o potencial passa de $2c-c=c$ para $2(c+1)-2c=2$, ou seja, $\Delta\Phi=2-c$, e o custo amortizado é $(c+1)+(2-c)=3$. Logo $n$ inserções custam no máximo $3n$.

\resp{hb-contador} Comece em $2^{k-1}-1=0111\ldots1$. Um incremento leva a $1000\ldots0$, trocando $k$ bits, e um decremento volta ao estado anterior, trocando $k$ bits. Alternando as duas operações, cada uma custa $k$, e o total é $\Theta(mk)$. Com apenas incrementos, o potencial ``número de bits iguais a 1'' pagava os ``vai um'' porque os uns, uma vez zerados, só voltam um a um; com decrementos, o mesmo trabalho pode ser repetido indefinidamente sem que nada o pague.

\resp{hb-fila} \emph{Enfileirar}: empilhe na pilha de entrada. \emph{Desenfileirar}: se a pilha de saída estiver vazia, transfira todos os elementos da entrada para ela, invertendo a ordem; depois desempilhe da saída. Cada elemento é empilhado e desempilhado no máximo uma vez em cada pilha, de modo que $m$ operações custam $O(m)$. Com o potencial $\Phi=2\cdot|\text{entrada}|$, enfileirar custa 3 amortizado e desenfileirar custa 1.

\resp{hb-forma} $100=1100100_2$: $\{B_2,B_5,B_6\}$. $101=1100101_2$: $\{B_0,B_2,B_5,B_6\}$, numa inserção sem links. De 100 a $104=1101000_2$: $\{B_3,B_5,B_6\}$. As quatro inserções ($100\to101\to102\to103\to104$) fazem $0,1,0,3$ links, respectivamente: o número de uns no final de $100$, $101$, $102$ e $103$.

\resp{hb-fibbin} Por indução. As árvores nascem com um único nó ($B_0$). Sem cortes, uma árvore só muda por \emph{link}, que na consolidação liga duas raízes de mesmo grau $k$; se ambas são árvores binomiais $B_k$ (não ordenadas), o resultado é $B_{k+1}$. Na remoção do máximo, os filhos da raiz de uma $B_k$ são raízes de $B_0,\dots,B_{k-1}$, também binomiais. Logo todas as árvores são binomiais; a diferença em relação ao heap binomial é que, antes da consolidação, pode haver várias árvores de mesma ordem.

\resp{hb-fibcaro} Insira $n$ elementos --- cada inserção só acrescenta uma raiz, e ficam $n$ raízes --- e depois remova o máximo: a consolidação percorre as $n-1$ raízes restantes e faz $\Theta(n)$ links, com custo real $\Theta(n)$. Não há contradição, porque o custo amortizado $O(\log n)$ vale para a \emph{sequência}: as $n$ inserções pagaram, cada uma com um crédito (o aumento de potencial $\Delta t=+1$), pelo trabalho da consolidação.

\resp{hb-binsim} Depois do 4: $\{4\}$. Depois do 9: as duas $B_0$ se juntam, com o 9 por cima: $\{9\to4\}$ (1 link). Depois do 2: $\{2,\ 9\to4\}$ (nenhum). Depois do 7: o 7 se junta ao 2, formando $7\to2$, que se junta a $9\to4$, com o 9 por cima: uma $B_2$ de raiz 9, cujos filhos são $7\to2$ e 4 (2 links). Depois do 5: $\{5,\ B_2\}$ (nenhum). Os links por inserção são $0,1,0,2,0$, exatamente os ``vai um'' das somas $0+1$, $1+1$, $10_2+1$, $11_2+1$ e $100_2+1$; o total é 3, e o heap final tem $5=101_2$ elementos.

\resp{hb-fibsim} As sete inserções criam sete raízes, e $\Phi=7$. O 9 sai, e a consolidação percorre $6,2,4,8,1,7$: o 2 colide com o 6, formando $6\to2$, de grau 1; o 4 fica em $A[0]$; o 8 colide com o 4, formando $8\to4$, que colide com $6\to2$: o 8 fica por cima, numa árvore de grau 2, com os filhos $6\to2$ e 4; o 1 vai para $A[0]$, e o 7 colide com ele, formando $7\to1$, de grau 1. Restam duas raízes, 8 (grau 2) e 7 (grau 1), depois de 4 links; o novo máximo é 8, e $\Phi$ cai de 7 para 2.
